\section{Discussion and falsification criteria}
At a high level, the project separates four questions: whether the action is coherent, whether a recurrent background exists, whether perturbations are healthy, and whether derived observables fit data. At a low level, each gate corresponds to a checkable artifact: conservation and constraint residuals with root-localized events; exact turning-state theorem audits and recurrence-candidate reports; Poincar\'e maps and Floquet multipliers; quadratic kinetic and gradient matrices; and release-native likelihoods with transfer functions and nuisance parameters. The machine-readable gate ledger prevents a later evidential claim from being marked passed while a prerequisite remains unresolved.

The present canonical baseline passes the numerical-consistency gate, but its status is now stronger than a finite-time non-detection of $H=0$. Because the configured canonical potential has $V_{\min}=3.5\Mpl^4$, every regular turning point must have $a_\star\le0.9258$, whereas the selected expanding trajectory starts at $a_0=1$. It is therefore analytically excluded from every future regular turnaround and from future cyclic recurrence. This is a trajectory-level falsification result for the configured baseline, not a numerical-horizon artifact. It does not exclude other parameter choices, particularly those with a lower or nonpositive potential floor, different initial conditions, or a modified action.

The analytic certificate also changes the most efficient Stage~1 search strategy. Expanding initial states that satisfy a rigorous future-turnaround exclusion should be classified before expensive long integrations. Numerical scans should concentrate on the remaining admissible region, retain the excluded points as explicit negative evidence, and use the exact event-state candidate protocol only where a turning sequence can actually occur. This improves statistical honesty as well as computational efficiency because the search does not preferentially retain visually interesting trajectories.

The DESI result has a similarly bounded interpretation. Reproducing the released flat-$\Lambda$CDM chain constraints demonstrates that the BAO vector, covariance ordering, ruler calibration, GetDist chain handling, and CAMB background machinery are functioning consistently. It does not test \SCPC{} until a calibrated \SCPC{} $H(z)$, $r_d$, and ultimately perturbation transfer functions have been derived. The July 2026 DESI Lyman-$\alpha$ full-shape/AP result is therefore retained as a separate comparator rather than being multiplied into the pinned 2025 likelihood with an invented cross-covariance. Its mock-validated BAO/AP information may become useful once an official joint likelihood and an admissible \SCPC{} observable map both exist, while the biased $f\sigma_8$ estimator reported by the validation program is explicitly excluded \cite{desi2026iv,desi2026lyavalidation}.

The next defensible calculations are therefore:
\begin{enumerate}
\item scan only the analytically admissible closed-FLRW parameter/initial-condition domain, while retaining exclusion certificates, failed runs, and domain-terminated runs as first-class outcomes;
\item refine outcome boundaries with predeclared deterministic or seeded low-discrepancy designs rather than visual hand-selection, and subject any putative cyclic point to the full multi-tolerance, Radau-versus-DOP853, exact-event candidate gate;
\item only for Stage~1 candidates that pass that gate, compute a Poincar\'e map, variational equations, monodromy matrix, Floquet multipliers, basin size, and tolerance/solver convergence;
\item derive the quadratic scalar and tensor perturbation actions and reject ghost, gradient-unstable, or strongly coupled regions;
\item only after those gates connect the background and perturbations to CLASS, CAMB, or another validated Boltzmann solver \cite{class2011,camb2000};
\item perform synthetic parameter recovery before attempting joint Planck/ACT, DESI, supernova, growth, lensing, or gravitational-wave inference.
\end{enumerate}

The canonical model family is rejected as a cyclic theory if no Stage~1 candidate survives a declared, converged admissible domain, or if every surviving background fails the later homogeneous or perturbative stability gates. A proposed extension is rejected if it violates stress-energy conservation, lacks a controlled standard-model limit, contains a ghost or unacceptable gradient instability, or places its strong-coupling scale below the calculation. A spectral feature is accepted as physical only if it survives resolution, duration, cadence, window, transfer-function, foreground, and look-elsewhere tests. Repeated close returns in a single integration are not, by themselves, evidence for a stable limit cycle.
